\subsection{Mean computation and centering}
Before actual applying principal component analysis, the time delay embedding vectors must be centered by subtracting the mean. In the online scenarios, computing the mean raises the challenge that as new observations $k + 1$ continuously arrive, older observations become  irrelevant. The algorithm must therefore be capable to efficiently updating the mean as the sliding window advances.
Given a sliding window buffer of fixed capacity $W$, the running mean vector $\mu_t \in R^d$ across the $d$ embedding dimensions at iteration $t$ is calculated over the $N$ currently valid embedded vectors ($N \le W - (d - 1)\tau$):
\begin{equation}
    \boldsymbol{\mu}_t = \frac{1}{N} \sum_{i=1}^{N} \mathbf{x}_i
    \label{eq:tde_mean}
\end{equation}
When new scalar stream samples $x_{new}$ arrive, a new embedded vector is formed at the window tail. 
Once the window is initially full, the oldest embedded vector $x_{old}$ is discarded, allowing the running mean to update incrementally in $O(d)$ time:
\begin{equation}
\mu_{t+1} = \mu_t + \frac{x_{new} - x_{old}}{N}
\label{eq:mean_update}
\end{equation}
Once the running mean is established it can be assessed across different functions only then the time delay embedding vector is centered 
\begin{equation}
\tilde{\mathbf{z}}_t = \mathbf{z}_t - \boldsymbol{\mu} = 
\begin{bmatrix}
z_{t,1} - \mu \\
z_{t,2} - \mu \\
\vdots \\
z_{t,d} - \mu
\end{bmatrix}
\label{eq:centering_vector}
\end{equation}
Data centering translates the coordinate system origin to the sample mean vector $\mu_t$. This ensures that the resulting covariance matrix accurately captures variance relative to the center, rather than second-moment distortions introduced by a non-zero mean offset.